\clearpage
\section{Fitting-window and sample-count checks}
\label{sec:window-comparison}
The four active mixture windows are Eq.~\eqref{eq:common-fit-windows}.
They retain, per ordered pair, 1255 small-Ising, 598 small-boson,
260 large-Ising, and 317 large-boson observations. The corresponding
minimum training counts after omitting one chain size are 1250, 593,
255, and 312. Every active fit therefore exceeds forty observations,
including every omitted-size training set. All fitting and display
masks exclude $\ell=4,5$, while the raw files preserve those values.
Table~\ref{tab:scaling-sampling} gives exact sampled endpoints and
distinct-size counts. All five retained subsystem lengths occur in
each of the four active groups.

As a separate check closer to zero, we retain the earlier free-power
fits on $(0,0.001]$. Additional sizes within the validated maxima
restored their sample counts after excluding $\ell=4,5$; the later
dense scan added one further observation at $N=6000$, $\ell=6$.
After the fixed-window refinement, these diagnostic fits contain
57 Ising and 116 boson geometries per direction. They are not
the fits used in the active mixture figures or principal tables.

The count can be reconstructed without examining any entropy values.
For cutoff $t$ and the chain-size set $\mathcal N_{\rm CFT}$ of the
chosen CFT in Eq.~\eqref{eq:extended-grid}, define
\begin{equation}
 n_N(t)=\max\{0,\min[10,\lfloor tN\rfloor]-5\},\qquad
 m(t)=\sum_{N\in\mathcal N_{\rm CFT}}n_N(t).
 \label{eq:fit-sample-count}
\end{equation}
Here $\lfloor u\rfloor$ is the greatest integer no larger than $u$;
$n_N$ counts the admitted lengths starting at six. A size-holdout
training set contains $m(t)-n_N(t)$ points. Both $m(t)$ at each adopted cutoff and every
such training count are checked against the minimum of forty.

For Ising, $N\leq8192$ means that only $\ell=6,7,8$ can enter
the separate $0.001$ diagnostic fit: $\ell=9,10$ would require $N\geq9000,10000$. Their computed
values remain in the figures. For the boson, every subsystem size
$\ell=6,\ldots,10$ occurs among the selected points, with
$N\leq16384$. Adding sizes increases the sampling density within
these limits; it does not lower the smallest accessible fraction.
Points sharing a chain or subsystem size are deterministic samples,
not forty statistically independent measurements.
Different $(N,\ell)$ can have the same fraction; they remain distinct
lattice observations and receive equal weight in Eq.~\eqref{eq:objective}.

The denser boson data also support the trial cutoff $0.00075$, with
69 points and at least 64 after omitting one size. The boson cutoff
$0.000625$ contains 41 points but can leave only 36 in a training
fold, so it is excluded. The remaining smaller trial windows fail
the same forty-point requirement. The free-fit tables below give
the qualifying smaller-window diagnostics in a compact form.
Each row estimates $Cz^r$ with both parameters free using
Eq.~\eqref{eq:objective}; theory enters only the subsequent comparison.

The coefficient and exponent are correlated, so a small exponent change
can produce a larger coefficient change when extrapolated to unit
fraction. Furthermore, a fixed range $\ell=6,\ldots,10$ keeps the
lattice resolution fixed as $N$ increases. Stable exponent recovery
does not require that every finite-$\ell$ amplitude equal its continuum
value. These are direct fits and comparisons; no extra lattice-effect
power law is assumed.
\begingroup\footnotesize\setlength{\tabcolsep}{4pt}
\begin{longtable}{llrrrrr}
\caption{Sample-count audit for the active mixture fits. Counts apply to each of the six directions. Every full fit and every omitted-size training set contains at least 40 distinct $(N,\ell)$ geometries. The fraction range is the actual sampled range within the corresponding window of Eq.~\eqref{eq:common-fit-windows}; active mixture figures display fractions through $0.2$. The number of sizes counts distinct retained $N$.}
\label{tab:scaling-sampling}\\\toprule
Interval & CFT & Samples & Sizes & Min. CV training & Smallest fraction & Largest fraction\\\midrule
\endfirsthead
\toprule
Interval & CFT & Samples & Sizes & Min. CV training & Smallest fraction & Largest fraction\\\midrule
\endhead
\bottomrule\endfoot
small & ising & 1255 & 278 & 1250 & $7.324\!\times\!10^{-4}$ & 0.02\\
small & boson & 598 & 124 & 593 & $3.662\!\times\!10^{-4}$ & 0.1\\
large & ising & 260 & 67 & 255 & $7.324\!\times\!10^{-4}$ & 0.002\\
large & boson & 317 & 68 & 312 & $3.662\!\times\!10^{-4}$ & 0.0049779\\
\end{longtable}

\begin{longtable}{lrrrrrr}
\caption{Free power-law fits on small-interval windows meeting the minimum of 40 observations even after omitting any one size. Every row estimates $C,r$ in $Cz^r$, with $z=x,r=\alpha$ for small intervals and $z=y,r=\beta$ for large intervals. Coefficients include all powers of $\pi$; all computed samples have $\ell=6,\ldots,10$. The cutoff is an upper bound, $m$ is the retained count, and the objective is Eq.~\eqref{eq:objective}. Only qualifying cutoffs are shown. These are a separate asymptotic-range check, not the adopted fits in Eq.~\eqref{eq:common-fit-windows}.}
\label{tab:ising-small-windows}\\\toprule
Direction & Cutoff & $m$ & $C$ & $r$ & $\eta_c$ [\%] & $\eta_r$ [\%]\\\midrule
\endfirsthead
\toprule
Direction & Cutoff & $m$ & $C$ & $r$ & $\eta_c$ [\%] & $\eta_r$ [\%]\\\midrule
\endhead
\bottomrule\endfoot
$I\mid \sigma$ & 0.001 & 57 & 0.20452 & 1.9992 & -0.53234 & -0.041954\\
$\sigma\mid I$ & 0.001 & 57 & 0.20453 & 1.9992 & -0.52997 & -0.04176\\
$I\mid \varepsilon$ & 0.001 & 57 & 14.577 & 4.018 & 12.232 & 0.45071\\
$\varepsilon\mid I$ & 0.001 & 57 & 14.574 & 4.018 & 12.212 & 0.45011\\
$\sigma\mid \varepsilon$ & 0.001 & 57 & 0.20476 & 1.9993 & -0.41875 & -0.033779\\
$\varepsilon\mid \sigma$ & 0.001 & 57 & 0.20471 & 1.9993 & -0.43938 & -0.035187\\
\end{longtable}

\begin{longtable}{lrrrrrr}
\caption{Free power-law fits on small-interval windows meeting the minimum of 40 observations even after omitting any one size. Every row estimates $C,r$ in $Cz^r$, with $z=x,r=\alpha$ for small intervals and $z=y,r=\beta$ for large intervals. Coefficients include all powers of $\pi$; all computed samples have $\ell=6,\ldots,10$. The cutoff is an upper bound, $m$ is the retained count, and the objective is Eq.~\eqref{eq:objective}. Only qualifying cutoffs are shown. These are a separate asymptotic-range check, not the adopted fits in Eq.~\eqref{eq:common-fit-windows}.}
\label{tab:boson-small-windows}\\\toprule
Direction & Cutoff & $m$ & $C$ & $r$ & $\eta_c$ [\%] & $\eta_r$ [\%]\\\midrule
\endfirsthead
\toprule
Direction & Cutoff & $m$ & $C$ & $r$ & $\eta_c$ [\%] & $\eta_r$ [\%]\\\midrule
\endhead
\bottomrule\endfoot
$00\mid ss$ & 0.001 & 116 & 0.82876 & 2.0087 & 0.76489 & 0.43312\\
$00\mid ss$ & $7.500\!\times\!10^{-4}$ & 69 & 0.89103 & 2.0185 & 8.336 & 0.92557\\
$ss\mid 00$ & 0.001 & 116 & 0.82876 & 2.0087 & 0.76489 & 0.43312\\
$ss\mid 00$ & $7.500\!\times\!10^{-4}$ & 69 & 0.89103 & 2.0185 & 8.336 & 0.92557\\
$00\mid 3s,s$ & 0.001 & 116 & 4.1499 & 2.0064 & 0.91377 & 0.32068\\
$00\mid 3s,s$ & $7.500\!\times\!10^{-4}$ & 69 & 4.3823 & 2.0138 & 6.565 & 0.69105\\
$3s,s\mid 00$ & 0.001 & 116 & 4.1499 & 2.0064 & 0.91387 & 0.32069\\
$3s,s\mid 00$ & $7.500\!\times\!10^{-4}$ & 69 & 4.3823 & 2.0138 & 6.5651 & 0.69105\\
$ss\mid 3s,s$ & 0.001 & 116 & 1.6645 & 2.0032 & 1.19 & 0.15955\\
$ss\mid 3s,s$ & $7.500\!\times\!10^{-4}$ & 69 & 1.713 & 2.0071 & 4.14 & 0.35488\\
$3s,s\mid ss$ & 0.001 & 116 & 1.6645 & 2.0032 & 1.1902 & 0.15956\\
$3s,s\mid ss$ & $7.500\!\times\!10^{-4}$ & 69 & 1.713 & 2.0071 & 4.1401 & 0.35488\\
\end{longtable}

\begin{longtable}{lrrrrrr}
\caption{Free power-law fits on large-interval windows meeting the minimum of 40 observations even after omitting any one size. Every row estimates $C,r$ in $Cz^r$, with $z=x,r=\alpha$ for small intervals and $z=y,r=\beta$ for large intervals. Coefficients include all powers of $\pi$; all computed samples have $\ell=6,\ldots,10$. The cutoff is an upper bound, $m$ is the retained count, and the objective is Eq.~\eqref{eq:objective}. Only qualifying cutoffs are shown. These are a separate asymptotic-range check, not the adopted fits in Eq.~\eqref{eq:common-fit-windows}.}
\label{tab:ising-large-windows}\\\toprule
Direction & Cutoff & $m$ & $C$ & $r$ & $\eta_c$ [\%] & $\eta_r$ [\%]\\\midrule
\endfirsthead
\toprule
Direction & Cutoff & $m$ & $C$ & $r$ & $\eta_c$ [\%] & $\eta_r$ [\%]\\\midrule
\endhead
\bottomrule\endfoot
$I\mid \sigma$ & 0.001 & 57 & 0.7149 & 0.26364 & 15.372 & 5.4547\\
$\sigma\mid I$ & 0.001 & 57 & 0.7149 & 0.26364 & 15.371 & 5.4543\\
$I\mid \varepsilon$ & 0.001 & 57 & 3.2493 & 1.9983 & -1.2326 & -0.086132\\
$\varepsilon\mid I$ & 0.001 & 57 & 3.2954 & 2 & 0.16796 & 0.0019731\\
$\sigma\mid \varepsilon$ & 0.001 & 57 & 0.15835 & 0.25189 & 2.2203 & 0.75771\\
$\varepsilon\mid \sigma$ & 0.001 & 57 & 0.16243 & 0.25504 & 4.8522 & 2.0148\\
\end{longtable}

\begin{longtable}{lrrrrrr}
\caption{Free power-law fits on large-interval windows meeting the minimum of 40 observations even after omitting any one size. Every row estimates $C,r$ in $Cz^r$, with $z=x,r=\alpha$ for small intervals and $z=y,r=\beta$ for large intervals. Coefficients include all powers of $\pi$; all computed samples have $\ell=6,\ldots,10$. The cutoff is an upper bound, $m$ is the retained count, and the objective is Eq.~\eqref{eq:objective}. Only qualifying cutoffs are shown. These are a separate asymptotic-range check, not the adopted fits in Eq.~\eqref{eq:common-fit-windows}.}
\label{tab:boson-large-windows}\\\toprule
Direction & Cutoff & $m$ & $C$ & $r$ & $\eta_c$ [\%] & $\eta_r$ [\%]\\\midrule
\endfirsthead
\toprule
Direction & Cutoff & $m$ & $C$ & $r$ & $\eta_c$ [\%] & $\eta_r$ [\%]\\\midrule
\endhead
\bottomrule\endfoot
$00\mid ss$ & 0.001 & 116 & 1.2264 & 0.99587 & -0.59355 & -0.41329\\
$00\mid ss$ & $7.500\!\times\!10^{-4}$ & 69 & 1.2068 & 0.99371 & -2.1781 & -0.62939\\
$ss\mid 00$ & 0.001 & 116 & 1.2264 & 0.99587 & -0.59482 & -0.41346\\
$ss\mid 00$ & $7.500\!\times\!10^{-4}$ & 69 & 1.2068 & 0.9937 & -2.184 & -0.63019\\
$00\mid 3s,s$ & 0.001 & 116 & 70.278 & 4.9903 & -6.4314 & -0.19312\\
$00\mid 3s,s$ & $7.500\!\times\!10^{-4}$ & 69 & 78.936 & 5.0066 & 5.096 & 0.13181\\
$3s,s\mid 00$ & 0.001 & 116 & 70.201 & 4.9902 & -6.534 & -0.19593\\
$3s,s\mid 00$ & $7.500\!\times\!10^{-4}$ & 69 & 78.905 & 5.0065 & 5.0553 & 0.13098\\
$ss\mid 3s,s$ & 0.001 & 116 & 3.2701 & 2.001 & -0.60142 & 0.050917\\
$ss\mid 3s,s$ & $7.500\!\times\!10^{-4}$ & 69 & 3.376 & 2.0054 & 2.6196 & 0.26766\\
$3s,s\mid ss$ & 0.001 & 116 & 3.2688 & 2.001 & -0.64071 & 0.048605\\
$3s,s\mid ss$ & $7.500\!\times\!10^{-4}$ & 69 & 3.3753 & 2.0053 & 2.5966 & 0.26651\\
\end{longtable}

\endgroup
